\paragraph{A∞-模。}\label{sub:equivariant-a-infinity-modules}
与上文对\trTwistedComplex{}的讨论相平行，我们也可以引入 \trAinfty-模的一个%
\trEquivariant{}版本。$\scr{A}$ 上的一个\textbf{\trEquivariant{} $A_\infty$-模}
（equivariant $A_\infty$-module）是一个 \trAinfty-模 $\scr{M}$，再配上一个%
\trModule\trHomomorphism{}%
$\epsilon_{\scr{M}}\colon\scr{M}\to\epsilon_{\scr{A}}^*\scr{M}$，其阶 $\ge 2$ 的各项为
零，并且 $\epsilon_{\scr{M}}$ 与 $\epsilon_{\scr{A}}^*(\epsilon_{\scr{M}})$ 的%
\trComposition{}无论按哪个次序都是恒等\trModule{}\trEndomorphism。具体地说，这份附加
数据由一族映射 $\epsilon_{\scr{M}}\colon\scr{M}(X)\to\scr{M}(\epsilon_{\scr{A}}(X))$
组成，对 $X$ 与 $\epsilon_{\scr{A}}(X)$ 的这两个映射互逆，并且它们与 $\epsilon_{\scr{A}}$
\trCompatible，其意义是
\begin{equation*}
\medmuskip=1.5mu\relax\thickmuskip=3mu\relax
\mu^d_{\scr{M}}\bigl(\epsilon_{\scr{M}}(b),\epsilon_{\scr{A}}(a_{d-1}),\ldots,
\epsilon_{\scr{A}}(a_1)\bigr)
=\epsilon_{\scr{M}}\bigl(\mu^d_{\scr{M}}(b,a_{d-1},\ldots,a_1)\bigr).
\end{equation*}
\trMorphism{}及其\trComposition{}均按非\trEquivariant{}的情形定义，并把所得的%
\trACategory{}记作 $\scr{Q}=\mathit{mod}^{\mathrm{eq}}(\scr{A})$。其上带有一个诱导的
朴素 $\ZZ/2$-作用，此作用在\trObject{}上是平凡的；也就是说，对
$t\in\hom_{\scr{Q}}(\scr{M}_0,\scr{M}_1)$，令
\begin{equation*}
\medmuskip=1.5mu\relax\thickmuskip=3mu\relax
(\epsilon_{\scr{Q}}t)^d(b,a_{d-1},\ldots,a_1)
=\epsilon_{\scr{M}_1}\bigl(t^d(\epsilon_{\scr{M}_0}(b),\epsilon_{\scr{A}}(a_{d-1}),
\ldots,\epsilon_{\scr{A}}(a_1))\bigr).
\end{equation*}

若 $Y$ 是 $\scr{A}$ 的一个 $\epsilon_{\scr{A}}$-不动\trObject，则与之相伴的%
\trYoneda{}\trModule{} $\scr{Y}$ 自带一个自然的\trEquivariant{}结构：在
$\scr{Y}(X)=\hom_{\scr{A}}(X,Y)$ 上取 $\epsilon_{\scr{Y}}=\epsilon_{\scr{A}}$。这定义了
一个\trCohomology\trFullyFaithful{}的 \trAFunctor{}%
$\scr{J}^{\mathrm{eq}}_{\scr{A}^{\mathrm{fix}}}\colon
\scr{A}^{\mathrm{fix}}\to\scr{Q}$，即\textbf{\trEquivariant\trYonedaEmbedding}
（equivariant Yoneda embedding），它与诸 $\ZZ/2$-作用严格交换（借助%
\cref{thm:tw-yoneda-combination-cff} 的一个\trEquivariant{}版本，还能把它扩充成一个%
\trCohomology\trFullyFaithful{}的\trFunctor{}%
$Tw^{\mathrm{eq}}\scr{A}\to\scr{Q}$）。另一方面，若 $\scr{M}$ 是一个\trEquivariant{}
\trAinfty-模，则只考虑诸空间 $\scr{M}(X)$ 在 $\epsilon_{\scr{A}}$-不动的 $X$ 处的
$\ZZ/2$-不变部分，就得到 $\scr{A}^{\mathrm{inv}}$ 上的一个 \trAinfty-模
$\scr{M}^{\mathrm{inv}}$。这连同对 $\ZZ/2$-不变\trNaturalTransformation{}的相应限制
一起，定义了\textbf{不动部分 \trAFunctor}（fixed part $A_\infty$-functor）
$\scr{Q}^{\mathrm{inv}}\to\mathit{mod}(\scr{A}^{\mathrm{inv}})$。最后，还可以干脆忘掉
诸映射 $\epsilon_{\scr{M}}$，这就给出一个\trFullyFaithful{}的 \trAFunctor{}%
$\scr{Q}\to\mathit{mod}(\scr{A})$。

给定一个带 $\ZZ/2$-作用的\trVectorSpace{}\trChainComplex{} $Z$ 以及一个%
\trEquivariant{}\trModule{} $\scr{M}$，\trTensorProduct{} $Z\otimes\scr{M}$ 仍是%
\trEquivariant{}的。其次，给定两个\trEquivariant{}\trModule{}之间的一个 $\ZZ/2$-不变%
\trHomomorphism，\eqref{eq:module-cone-definition} 意义下的\trCone{}仍是一个%
\trEquivariant{}\trModule。最后，若 $Y$ 是一个 $\epsilon_{\scr{A}}$-不动\trObject{}且
$\scr{M}$ 是一个\trEquivariant{}\trModule，则扭转\trModule{} $\scr{T}_Y\scr{M}$ 自然是%
\trEquivariant{}的，并且在 $H(\scr{Q}^{\mathrm{inv}})$ 中有\trExactTriangle{}
\eqref{eq:twist-exact-triangle}。

\vkthmsync{4}
\begin{lemma}\label{thm:equivariant-twist-yoneda-invariant}
设 $Y_0$、$Y_1$、$Y_2$ 都是 $\epsilon_{\scr{A}}$-不动\trObject，而且
$\Hom_{H(\scr{A})}(Y_0,Y_1)$ 上的 $\ZZ/2$-作用是平凡的。假设在 $H^0(\scr{Q})$ 中有
\begin{equation}\label{eq:equivariant-twist-yoneda}
\scr{T}_{Y_0}(\scr{Y}_1)\cong\scr{Y}_2
\end{equation}
且此\trIsomorphism{}属于此\trCategory{}中\trMorphism{}群的 $\ZZ/2$-不变部分。那么，
一个平行的\trIsomorphism{} \eqref{eq:equivariant-twist-yoneda} 在
$H^0(\mathit{mod}(\scr{A}^{\mathrm{inv}}))$ 中也成立，其中诸\trYoneda{}\trModule{}与
扭转都是对 $\scr{A}^{\mathrm{inv}}$ 而取的。最后，若 $\scr{A}$ 是\trCfinite{}的，则在
$D(\scr{A}^{\mathrm{inv}})$ 中有 $T_{Y_0}(Y_1)\cong Y_2$。
\end{lemma}

\begin{proof}
为了把这几个\trCategory{}彼此区分清楚，我们把\trEquivariant{}\trYoneda{} \trAinfty-模
记作 $\scr{J}^{\mathrm{eq}}_{\scr{A}}(Y)$，而把它们在
$\mathit{mod}(\scr{A}^{\mathrm{inv}})$ 中的对应物记作
$\scr{J}_{\scr{A}^{\mathrm{inv}}}(Y)$。施用不动部分\trFunctor{}%
$\scr{Q}^{\mathrm{inv}}\to\mathit{mod}(\scr{A}^{\mathrm{inv}})$，就把
$\scr{J}^{\mathrm{eq}}_{\scr{A}}(Y_2)$ 送到 $\scr{J}_{\scr{A}^{\mathrm{inv}}}(Y_2)$。
另一方面，$\scr{T}_{Y_0}\bigl(\scr{J}^{\mathrm{eq}}_{\scr{A}}(Y_1)\bigr)$ 的像是
$\scr{A}^{\mathrm{inv}}$ 上的如下 \trAinfty-模 $\scr{C}$：
\begin{equation*}
\medmuskip=1.5mu\relax\thickmuskip=3mu\relax
\scr{C}(X)=\bigl(\hom_{\scr{A}}(Y_0,Y_1)\otimes\hom_{\scr{A}}(X,Y_0)[1]\bigr)^{\ZZ/2}
\oplus\hom_{\scr{A}}(X,Y_1)^{\ZZ/2},
\end{equation*}
并配有\eqref{eq:abstract-twist-module} 中那些通常的\trStructureMap。由于
$\hom_{\scr{A}}(Y_0,Y_1)$ 的反不变部分是\trAcyclic{}的，把第一个\trDirectSummand{}换成
$\hom_{\scr{A}}(Y_0,Y_1)^{\ZZ/2}\otimes\hom_{\scr{A}}(X,Y_0)^{\ZZ/2}$ 便得到一个与之%
\trQuasiisomorphic{}的子模；由\cref{thm:module-composition-quasiisomorphism}，它在
$\mathit{mod}(\scr{A}^{\mathrm{inv}})$ 中是一个\trIsomorphic{}的\trObject。而按定义，
这个子模恰恰就是 $\scr{T}_{Y_0}\bigl(\scr{J}_{\scr{A}^{\mathrm{inv}}}(Y_1)\bigr)$。
最后的那个论断则由\eqref{eq:twist-notions-compatible} 立即得出。
\end{proof}

\begin{remark}\label{thm:naive-group-actions-remarks}
我们以几条略嫌含混的一般性评论来收尾。对 \trACategory{}上的群作用作朴素处理，这种做法
与\trHomotopy{}代数结构的哲学并不相合；尽管如此，对有限群以及特征为 $0$（或特征 $p$
适当大，$p\gg 0$）的\trCoefficients{}域 $\KK$，它仍然工作得相当好。我们将在 (10b) 中
重新考虑一般的情形；不过，那里提出的定义虽比这里所用的更为精致，却也至多是一个初步的
近似。

我们对不动\trObject{}的用法也有些朴素。在一个恰当的\trCategory{}论表述中，人们应当
转而考虑“\trEquivariant{}\trObject”（equivariant objects），它们在群作用之下只在
相差若干特选\trIsomorphism{}的意义下不动。这一点的\trCohomology{}层次版本稍后
（在 (14b) 中）会简要出现，但要定义一个恰当的上链层次概念却并非直截了当。形式地说，
这种做法的好处本应在于：人们可以把 $Tw^{\mathrm{eq}}\scr{A}$ 与
$\mathit{mod}^{\mathrm{eq}}(\scr{A})$ 都定义为由不动\trObject{}构成的子\trCategory{}——分别对
$Tw\scr{A}$ 与 $\mathit{mod}(\scr{A})$ 的诱导作用而言。不过，我们在这里作出的这些特设
定义，对我们的目的来说已经足够了。
\end{remark}
